% --- Archon physics formalization source begin ---
% archon:physics
% archon:covers IPhO2026Problems/problem_IPhO_2026_2_B_3.lean
% archon:source-report reports/ipho_2026/problem_IPhO_2026_2_B_3.source.json
% archon:problem-id IPhO_2026_2
% archon:part-id B.3
% archon:previous-part IPhO_2026_2_B_1
% archon:previous-part-policy natural_language_prerequisite_only
% archon:previous-part IPhO_2026_2_B_2
% archon:previous-part-policy natural_language_prerequisite_only

\chapter{Physics problem IPhO\_2026\_2\_B\_3}
\label{ch:IPhO2026Problems_problem_IPhO_2026_2_B_3}

\paragraph{Problem source.}
A half-cylindrical mirror of radius R illuminates a fully absorbing cylindrical
container of radius a.  Their axes are parallel, and the container center lies
R/2 from the mirror center on the symmetry plane.  Uniform parallel sunlight
arrives along the optical axis.  Any ray absorbed by the container reflects at
most once.  Let theta\_max be the largest incidence angle on the mirror among
rays that strike the container, and let P\_0 be the power the cylinder would
receive without the mirror.  See Figure 2f.

Current subquestion:
For R = 1.0 m, find a such that P = 5*P\_0, and report it in cm.

\paragraph{Current subquestion.}
For R = 1.0 m, find a such that P = 5*P\_0, and report it in cm.

\paragraph{Recorded answer/context.}
cos(theta\_max) = 4/5 and a = 0.12 m = 12 cm.

\paragraph{Figure/image path.}
/root/proposal\_for\_physic/science-mango/ipho\_2026\_source/image/T2\_page-3.png

\paragraph{Reusable previous-part conclusions.}
\begin{itemize}
\item Source B.1. Question: Given a = alpha*sin(theta\_max) + beta*sin(2*theta\_max), determine alpha and beta in terms of R. Reusable conclusions: alpha = R and beta = -R/2. Policy: natural\_language\_prerequisite\_only; do\_not\_import\_Lean\_output
\item Source B.2. Question: Express the power ratio P/P\_0 in terms of theta\_max. Reusable conclusions: P/P\_0 = 1/(1 - cos(theta\_max)). Policy: natural\_language\_prerequisite\_only; do\_not\_import\_Lean\_output
\end{itemize}

\paragraph{Formalization target.}
create a compiling Lean file with sorry bodies at `IPhO2026Problems/problem\_IPhO\_2026\_2\_B\_3.lean`. The Lean declarations must preserve the physical quantities, dimensions or dimensional roles, figure labels, governing-law hypotheses, and final relation expressed by this problem.
Use Mathlib/Physlib names found through LeanExplore where available. If a physics API is missing, introduce faithful local abstractions rather than scalar placeholder aliases.

\begin{definition}[PhysicalLength]
  \label{decl:physics:IPhO_2026_2_B_3:PhysicalLength}
  \lean{IPhO2026Problems.IPhO2026_2_B_3.PhysicalLength}
  A physical length, represented independently of the choice of units.
\end{definition}
\begin{proof}
  This is the defining declaration; unfolding it gives the stated typed data or relation.
\end{proof}

\begin{definition}[OpticalPower]
  \label{decl:physics:IPhO_2026_2_B_3:OpticalPower}
  \lean{IPhO2026Problems.IPhO2026_2_B_3.OpticalPower}
  A physical power.
\end{definition}
\begin{proof}
  This is the defining declaration; unfolding it gives the stated typed data or relation.
\end{proof}

\begin{definition}[centimeterUnits]
  \label{decl:physics:IPhO_2026_2_B_3:centimeterUnits}
  \lean{IPhO2026Problems.IPhO2026_2_B_3.centimeterUnits}
  Unit choices obtained from SI by measuring length in centimetres.
\end{definition}
\begin{proof}
  This is the defining declaration; unfolding it gives the stated typed data or relation.
\end{proof}

\begin{definition}[lengthInMeters]
  \label{decl:physics:IPhO_2026_2_B_3:lengthInMeters}
  \lean{IPhO2026Problems.IPhO2026_2_B_3.lengthInMeters}
  \uses{decl:physics:IPhO_2026_2_B_3:PhysicalLength}
  The numerical readout of a physical length in metres.
\end{definition}
\begin{proof}
  This is the defining declaration; unfolding it gives the stated typed data or relation.
\end{proof}

\begin{definition}[lengthInCentimeters]
  \label{decl:physics:IPhO_2026_2_B_3:lengthInCentimeters}
  \lean{IPhO2026Problems.IPhO2026_2_B_3.lengthInCentimeters}
  \uses{decl:physics:IPhO_2026_2_B_3:PhysicalLength, decl:physics:IPhO_2026_2_B_3:centimeterUnits}
  The numerical readout of a physical length in centimetres.
\end{definition}
\begin{proof}
  This is the defining declaration; unfolding it gives the stated typed data or relation.
\end{proof}

\begin{definition}[powerInSI]
  \label{decl:physics:IPhO_2026_2_B_3:powerInSI}
  \lean{IPhO2026Problems.IPhO2026_2_B_3.powerInSI}
  \uses{decl:physics:IPhO_2026_2_B_3:OpticalPower}
  The numerical SI readout of an optical power.
\end{definition}
\begin{proof}
  This is the defining declaration; unfolding it gives the stated typed data or relation.
\end{proof}

\begin{definition}[HalfCylindricalMirror]
  \label{decl:physics:IPhO_2026_2_B_3:HalfCylindricalMirror}
  \lean{IPhO2026Problems.IPhO2026_2_B_3.HalfCylindricalMirror}
  \uses{decl:physics:IPhO_2026_2_B_3:PhysicalLength}
  The half-cylindrical mirror appearing in Figure 2f.
\end{definition}
\begin{proof}
  This is the defining declaration; unfolding it gives the stated typed data or relation.
\end{proof}

\begin{definition}[FullyAbsorbingCylinder]
  \label{decl:physics:IPhO_2026_2_B_3:FullyAbsorbingCylinder}
  \lean{IPhO2026Problems.IPhO2026_2_B_3.FullyAbsorbingCylinder}
  \uses{decl:physics:IPhO_2026_2_B_3:PhysicalLength}
  The fully absorbing cylindrical container appearing in Figure 2f.
\end{definition}
\begin{proof}
  This is the defining declaration; unfolding it gives the stated typed data or relation.
\end{proof}

\begin{definition}[SunlightBeam]
  \label{decl:physics:IPhO_2026_2_B_3:SunlightBeam}
  \lean{IPhO2026Problems.IPhO2026_2_B_3.SunlightBeam}
  The uniform parallel sunlight incident along the optical axis.
\end{definition}
\begin{proof}
  This is the defining declaration; unfolding it gives the stated typed data or relation.
\end{proof}

\begin{definition}[Figure2fSetup]
  \label{decl:physics:IPhO_2026_2_B_3:Figure2fSetup}
  \lean{IPhO2026Problems.IPhO2026_2_B_3.Figure2fSetup}
  \uses{decl:physics:IPhO_2026_2_B_3:PhysicalLength, decl:physics:IPhO_2026_2_B_3:HalfCylindricalMirror, decl:physics:IPhO_2026_2_B_3:FullyAbsorbingCylinder, decl:physics:IPhO_2026_2_B_3:SunlightBeam}
  The physical objects and center-to-center separation shown in Figure 2f.
\end{definition}
\begin{proof}
  This is the defining declaration; unfolding it gives the stated typed data or relation.
\end{proof}

\begin{theorem}[Physics formalization target]
\label{thm:physics:IPhO_2026_2_B_3:target}
\lean{IPhO2026Problems.IPhO2026_2_B_3.radius_for_fivefold_power}
\uses{decl:physics:IPhO_2026_2_B_3:PhysicalLength, decl:physics:IPhO_2026_2_B_3:OpticalPower, decl:physics:IPhO_2026_2_B_3:centimeterUnits, decl:physics:IPhO_2026_2_B_3:lengthInMeters, decl:physics:IPhO_2026_2_B_3:lengthInCentimeters, decl:physics:IPhO_2026_2_B_3:powerInSI, decl:physics:IPhO_2026_2_B_3:HalfCylindricalMirror, decl:physics:IPhO_2026_2_B_3:FullyAbsorbingCylinder, decl:physics:IPhO_2026_2_B_3:SunlightBeam, decl:physics:IPhO_2026_2_B_3:Figure2fSetup}
The assigned autoformalize agent should translate this physics problem into Lean declarations in the covered file, with theorem and lemma proof bodies written as `by sorry`.
\end{theorem}
\begin{proof}
This is an autoformalization task, not a proof task. Produce statements that can later be proved by the physics prover without weakening the physical model.
\end{proof}
% --- Archon physics formalization source end ---
